\begin{frame}{例子：算一遍 KV 容量为什么能立刻改变系统判断}
  \Core{KV 课程里最应该让学生亲手做的，是容量算术，因为它直接决定你如何理解并发上限。}
  \begin{block}{输入更长}
    单请求 KV 规模线性增长，prefill 结束后就把一大块状态留在系统里。
  \end{block}
  \pause
  \begin{block}{并发更高}
    多请求同时驻留时，容量压力不是加法直觉那么简单，还会叠加块管理与碎片问题。
  \end{block}
  \pause
  \begin{block}{课堂结论}
    只要状态规模算错，后续对调度、throughput 和 tail latency 的解释就会整体偏掉。
  \end{block}
\end{frame}

\begin{frame}{例子：为什么连续大块分配很快会失败}
  \Core{传统“给每个请求预留一整段连续显存”的想法，在动态服务下几乎注定碰壁。}
  \begin{block}{预留过大}
    会浪费大量空闲容量，使理论并发远高于实际并发。
  \end{block}
  \pause
  \begin{block}{预留过小}
    一旦上下文稍长，请求就会频繁扩容甚至失败。
  \end{block}
  \pause
  \begin{block}{PagedAttention 的意义}
    把逻辑连续变成物理分块后，系统终于可以更诚实地管理状态生命周期。
  \end{block}
\end{frame}

\begin{frame}{代码例子：BlockManager 里最值得学生先看的对象}
  \Core{KV 讲义里的代码例子，应该从对象关系入手，而不是从论文术语入手。}
  \begin{block}{核心对象}
    {\ttfamily\small
    \detokenize{Block(block_id, ref_count, hash_value, token_ids)}\par
    \detokenize{free_block_ids}\par
    \detokenize{used_block_ids}\par
    \detokenize{hash_to_block_id}\par
    }
  \end{block}
  \pause
  \begin{block}{课堂应追问}
    这些对象分别体现了分配、占用、共享和回收的哪一面？
  \end{block}
  \pause
  \begin{block}{系统含义}
    读懂这些对象，才会知道 KV 管理讲的并不是“抽象缓存”，而是明确的数据结构与状态机。
  \end{block}
\end{frame}

\begin{frame}{代码例子：从 can\_allocate 到 allocate 的判断链}
  \Core{课堂最适合拿来解释的，不是整份代码，而是“能否分配 -> 如何分配”的判断链。}
  \begin{block}{代码片段}
    {\ttfamily\small
    \detokenize{num_required_blocks = seq.num_blocks - num_computed_blocks}\par
    \detokenize{return num_required_blocks <= len(self.free_block_ids)}\par
    \detokenize{seq.block_table.append(block_id)}\par
    }
  \end{block}
  \pause
  \begin{block}{课堂应追问}
    为什么“还剩多少 free blocks”本身就已经在决定系统能否扩大并发？
  \end{block}
  \pause
  \begin{block}{结论}
    KV 不是调度的附属品，容量判断本身就在塑造调度器的可选动作空间。
  \end{block}
\end{frame}

\begin{frame}{例子：前缀复用为什么比普通缓存更复杂}
  \Core{一旦系统尝试前缀共享，KV 就从“每请求私有状态”变成“可能跨请求共享状态”。}
  \begin{block}{收益直觉}
    相同前缀意味着重复 prefill 可被避免，潜在节省显著。
  \end{block}
  \pause
  \begin{block}{复杂性来源}
    共享需要元数据索引、引用计数和正确回收边界，否则收益可能变成一致性问题。
  \end{block}
  \pause
  \begin{block}{课程边界}
    课堂里可以先讲 observe-only 的潜在命中，再逐步进入真实复用机制。
  \end{block}
\end{frame}

\begin{frame}{课堂练习：把 KV 问题翻译成系统因果链}
  \Core{KV 这一讲不应该只留下术语，而应该留下至少一条学生自己能复述的因果链。}
  \begin{block}{练习一}
    用一句因果链解释：为什么长上下文会通过状态驻留拖慢短请求体验？
  \end{block}
  \pause
  \begin{block}{练习二}
    说明为什么分页/块管理的系统意义不只是“节省显存”。
  \end{block}
  \pause
  \begin{block}{练习三}
    如果你要为 BlockManager 加第一条观测日志，你会放在 can\_allocate、allocate 还是 deallocate？为什么？
  \end{block}
\end{frame}
